% Generated by `taskgen course examples` from crypto.lcg_period (seed 15).
% Do not edit: change examples.yaml and regenerate.
\begin{example}[проверка условий теоремы о максимальном периоде]\label{ex:lcg-full-period}
Имеет ли линейная конгруэнтная последовательность $x_{n+1} = (a x_n + c) \bmod m$ с параметрами \[a = 31,\quad c = 131,\quad m = 300\] максимальный период $m$ (при любом $x_0$)? Если нет, указать нарушенные условия теоремы о максимальном периоде.

\par\noindent\textit{Решение.} Теорема (Халл"--~Добелл; Кнут, теорема~A): период равен $m$ тогда и только тогда, когда 1)~$\gcd(c, m) = 1$; 2)~$p \mid a - 1$ для каждого простого $p \mid m$; 3)~$4 \mid a - 1$, если $4 \mid m$.\par Здесь $m = 2^{2} \cdot 3 \cdot 5^{2}$, $a - 1 = 2 \cdot 3 \cdot 5$.\par 1) $\gcd(131, 300) = 1$ --- выполнено.\par 2) $2 \mid 30$, $3 \mid 30$, $5 \mid 30$ --- выполнено.\par 3) $4 \mid m$, $a - 1 \equiv 2 \pmod 4$ --- \textbf{нарушено}.\[\text{период меньше $m$ (нарушены условия 3)}.\]
\end{example}
